% Supplementary Information of the ERL Letter (D-107): accuracy of the waterline elevation
% model and conformal coverage (H2), against the four independent height sources (ICESat-2
% ATL06 and ATL08, GLO-30, GEDI). D-086, D-088, D-091. Every number is a \metric{} macro.

\subsection*{Verdict}
H2, that the waterline elevation model has lower vertical error than the grids, is \metric{h2.decision.matched} under the pre-registered rule (matched variant; \metric{h2.decision.total} in the total variant), for both ICESat-2 products (table~\ref{tab:h2}). On \metric{h2.atl08.n_test_blocks} held-out test blocks of \metric{h2.block_size_km}~km with \metric{h2.atl08.n_test_points} ATL08 points, the model's NMAD is \metric{h2.atl08.dem.nmad_m}~m against \metric{erl.text.h2.atl08.product_nmad_min_m}--\metric{erl.text.h2.atl08.product_nmad_max_m}~m for the grids, but its root-mean-square error (RMSE) is \metric{h2.atl08.dem.rmse_m}~m against \metric{erl.text.h2.atl08.product_rmse_min_m}--\metric{erl.text.h2.atl08.product_rmse_max_m}~m; with ATL06 (\metric{h2.atl06.n_test_points} points) the values are \metric{h2.atl06.dem.nmad_m}~m against \metric{erl.text.h2.atl06.product_nmad_min_m}--\metric{erl.text.h2.atl06.product_nmad_max_m}~m (NMAD) and \metric{h2.atl06.dem.rmse_m}~m against \metric{erl.text.h2.atl06.product_rmse_min_m}--\metric{erl.text.h2.atl06.product_rmse_max_m}~m (RMSE). The paired block-bootstrap \metric{params.ci_pct}\% intervals of the differences include zero for \metric{erl.text.h2.atl08.diff_nmad.n_ci_include_zero} of \metric{erl.text.h1.n_products} grids in NMAD and \metric{erl.text.h2.atl08.diff_rmse.n_ci_include_zero} of \metric{erl.text.h1.n_products} in RMSE (ATL08). The brief's pre-registered fallback therefore applies: the audit of the Letter uses the observed extents, not the elevation model, and the model enters only the fused surface of the near-term analysis, above the record low.

\emph{Where it fails.} The model is good where the seabed is bare and the level is basin-wide (the interfluve, northeastern coast and open shelf) and fails in two places that decide the verdict (table~\ref{tab:ed_h2_sectors}). At the Volga delta front its ATL08 bias is \metric{h2.atl08.dem.volga.bias_m}~m, probably because flooded reeds that the vegetation rule missed are classified as water. In Komsomolets Bay its bias is \metric{h2.atl08.dem.komsomolets.bias_m}~m: the bay's water level departs from the basin level (setdown or partial isolation), which breaks the pairing of waterlines with the basin level. On vegetated seabed, a separate population (robustness check~9), the model's ATL08 NMAD is \metric{h2.veg.atl08.dem.nmad_m}~m (bias \metric{h2.veg.atl08.dem.bias_m}~m) against \metric{h2.veg.atl08.etopo2022.nmad_m}~m for ETOPO 2022 (figure~\ref{fig:ed_vegetation}).

\emph{Block size.} The verdict depends on the block size (robustness check~8; figure~\ref{fig:ed_blocksize}). With \metric{params.block_size_km}~km blocks (\metric{h2.sens25km.atl08.n_test_blocks} ATL08 test blocks) the rule is still not met: the model's NMAD advantage is significant against every grid, but the RMSE intervals include zero. With \metric{params.block_size_sens_lo_km}~km blocks (\metric{h2.sens12km.atl08.n_test_blocks} test blocks) H2 is supported for ATL08 in both variants, and for ATL06 in the total variant.

\emph{Which blocks are held out.} The verdict also depends on the draw of the blocks (robustness check~8). In \metric{robustness.check08.split_reassignment.n_seeds} random reassignments of the \metric{h2.block_size_km}~km blocks to training, calibration and test, with the same rule and ATL08, H2 is supported in \metric{robustness.check08.split_reassignment.atl08.n_seeds_h2_matched_supported} (matched variant; \metric{robustness.check08.split_reassignment.atl08.n_seeds_h2_total_supported} in the total variant). Across the draws the model's test-block NMAD is \metric{robustness.check08.split_reassignment.atl08.dem_nmad_min_m}--\metric{robustness.check08.split_reassignment.atl08.dem_nmad_max_m}~m (median \metric{erl.text.h2.split.atl08.nmad_median_m}~m) and its median coverage \metric{erl.text.h2.split.atl08.coverage_median_pct}\%; the registered draw has a higher NMAD than \metric{erl.text.h2.split.atl08.n_seeds_lower_nmad} of them and lower coverage than \metric{erl.text.h2.split.atl08.n_seeds_higher_coverage}, so it lies on the unfavourable side of the spread. These are post-hoc sensitivities that reuse test blocks. The registered verdict at the registered block size and draw stands, as registered, but it rests on \metric{h2.atl08.n_test_blocks} test blocks and on which blocks they are.

\subsection*{Conformal coverage}
The nominal \metric{erl.text.conformal_pct}\% intervals cover \metric{erl.text.pct.h2.atl06.coverage.overall}\% (ATL06) and \metric{erl.text.pct.h2.atl08.coverage.overall}\% (ATL08) of test heights. The Clopper--Pearson \metric{params.ci_pct}\% intervals (\metric{erl.text.pct.h2.atl08.coverage.overall_ci_low}--\metric{erl.text.pct.h2.atl08.coverage.overall_ci_high}\% for ATL08) exclude \metric{erl.text.conformal_pct}\%, so coverage is \metric{h2.atl08.coverage.overall_consistent} with the nominal level, overall and in every sector; they treat points as independent and are optimistic, and the test-block bootstrap intervals are \metric{erl.text.pct.h2.atl06.coverage.overall_boot_low}--\metric{erl.text.pct.h2.atl06.coverage.overall_boot_high}\% (ATL06) and \metric{erl.text.pct.h2.atl08.coverage.overall_boot_low}--\metric{erl.text.pct.h2.atl08.coverage.overall_boot_high}\% (ATL08). Coverage is near zero in Komsomolets Bay (\metric{erl.text.pct.h2.atl08.coverage.komsomolets}\%, ATL08), whose local bias no training or calibration block shares, so block-level exchangeability fails there; elsewhere the intervals over- or under-cover by sector. The interval method is reported as not achieving its nominal level; it was not tuned.

\subsection*{The other independent heights}
GLO-30 and GEDI, both pre-registered independent checks, were run after the ICESat-2 decision, with the same blocks, population and pipeline and their own training-block constants, and never pooled (table~\ref{tab:si_glo30_gedi}). They give the same picture: the model's NMAD is lower than most grids', its RMSE is not, and the rule is not met; it fails at the Volga delta front and in Komsomolets Bay against every source. Coverage is \metric{erl.text.pct.h2.glo30.coverage.overall}\% against GLO-30 and \metric{erl.text.pct.h2.gedi.coverage.overall}\% against GEDI, not consistent with \metric{erl.text.conformal_pct}\%.

GLO-30 is not independent of ETOPO 2022: in the emerged band ETOPO 2022 is filled from FABDEM, which is built from the Copernicus elevation model. \metric{h2.glo30.etopo_src.fabdem.share}\% of the GLO-30 test points lie in ETOPO 2022's FABDEM cells, where its NMAD against GLO-30 is \metric{h2.glo30.etopo_src.fabdem.etopo2022.nmad_m}~m (the model's, \metric{h2.glo30.etopo_src.fabdem.dem.nmad_m}~m); in its lake-depth cells, \metric{h2.glo30.etopo_src.gebco_lake.etopo2022.nmad_m}~m against the model's \metric{h2.glo30.etopo_src.gebco_lake.dem.nmad_m}~m. GEDI is independent of both. Against GEDI, ETOPO 2022's NMAD is \metric{h2.gedi.etopo_src.fabdem.etopo2022.nmad_m}~m in its FABDEM cells (\metric{h2.gedi.etopo_src.fabdem.share}\% of points; model \metric{h2.gedi.etopo_src.fabdem.dem.nmad_m}~m) and \metric{h2.gedi.etopo_src.gebco_lake.etopo2022.nmad_m}~m in its lake-depth cells (model \metric{h2.gedi.etopo_src.gebco_lake.dem.nmad_m}~m). ETOPO 2022's advantage over the other grids in the upper band is therefore real, GLO-30 overstates it, and it does not extend to the cells filled from lake depths. GLO-30 samples only seabed that was dry before December 2010, and most of its test points fall in two blocks on the northeastern coast.

\begin{table}[!htbp]
\centering
\small
\caption{\textbf{The waterline elevation model against GLO-30 and GEDI heights.} Held-out test blocks of \metric{h2.block_size_km}~km, bare emerged seabed, matched variant; errors are model minus reference (m). Each source has its own training-block constant $c$ (with its block-bootstrap standard deviation). $\Delta$, model minus grid, with the paired test-block bootstrap \metric{params.ci_pct}\% interval; negative favours the model. GEDI night: shots with the Sun below the horizon only (pre-registration Amendment~2).}
\label{tab:si_glo30_gedi}
\begin{tabular}{lccc}
\hline
 & GLO-30 & GEDI & GEDI, night \\
\hline
Test points (blocks) & \metric{h2.glo30.n_test_points} (\metric{h2.glo30.n_test_blocks}) & \metric{h2.gedi.n_test_points} (\metric{h2.gedi.n_test_blocks}) & \metric{h2.gedi_night.n_test_points} (\metric{h2.gedi_night.n_test_blocks}) \\
Constant $c$ (m) & \metric{h2.glo30.datum_constant_m} $\pm$ \metric{h2.glo30.datum_constant_sigma_m} & \metric{h2.gedi.datum_constant_m} $\pm$ \metric{h2.gedi.datum_constant_sigma_m} & as GEDI \\
Model bias / RMSE / NMAD & \metric{h2.glo30.dem.bias_m} / \metric{h2.glo30.dem.rmse_m} / \metric{h2.glo30.dem.nmad_m} & \metric{h2.gedi.dem.bias_m} / \metric{h2.gedi.dem.rmse_m} / \metric{h2.gedi.dem.nmad_m} & \metric{h2.gedi_night.dem.bias_m} / \metric{h2.gedi_night.dem.rmse_m} / \metric{h2.gedi_night.dem.nmad_m} \\
GEBCO\_2026 RMSE / NMAD & \metric{h2.glo30.gebco2026.rmse_m} / \metric{h2.glo30.gebco2026.nmad_m} & \metric{h2.gedi.gebco2026.rmse_m} / \metric{h2.gedi.gebco2026.nmad_m} & \metric{h2.gedi_night.gebco2026.rmse_m} / \metric{h2.gedi_night.gebco2026.nmad_m} \\
ETOPO 2022 RMSE / NMAD & \metric{h2.glo30.etopo2022.rmse_m} / \metric{h2.glo30.etopo2022.nmad_m} & \metric{h2.gedi.etopo2022.rmse_m} / \metric{h2.gedi.etopo2022.nmad_m} & \metric{h2.gedi_night.etopo2022.rmse_m} / \metric{h2.gedi_night.etopo2022.nmad_m} \\
$\Delta$NMAD, GEBCO\_2026 & \metric{h2.glo30.diff.nmad.gebco2026_m} [\metric{h2.glo30.diff.nmad.gebco2026_ci_low_m}, \metric{h2.glo30.diff.nmad.gebco2026_ci_high_m}] & \metric{h2.gedi.diff.nmad.gebco2026_m} [\metric{h2.gedi.diff.nmad.gebco2026_ci_low_m}, \metric{h2.gedi.diff.nmad.gebco2026_ci_high_m}] & \metric{h2.gedi_night.diff.nmad.gebco2026_m} [\metric{h2.gedi_night.diff.nmad.gebco2026_ci_low_m}, \metric{h2.gedi_night.diff.nmad.gebco2026_ci_high_m}] \\
$\Delta$RMSE, ETOPO 2022 & \metric{h2.glo30.diff.rmse.etopo2022_m} [\metric{h2.glo30.diff.rmse.etopo2022_ci_low_m}, \metric{h2.glo30.diff.rmse.etopo2022_ci_high_m}] & \metric{h2.gedi.diff.rmse.etopo2022_m} [\metric{h2.gedi.diff.rmse.etopo2022_ci_low_m}, \metric{h2.gedi.diff.rmse.etopo2022_ci_high_m}] & \metric{h2.gedi_night.diff.rmse.etopo2022_m} [\metric{h2.gedi_night.diff.rmse.etopo2022_ci_low_m}, \metric{h2.gedi_night.diff.rmse.etopo2022_ci_high_m}] \\
Coverage of \metric{erl.text.conformal_pct}\% intervals (\%) & \metric{erl.text.pct.h2.glo30.coverage.overall} & \metric{erl.text.pct.h2.gedi.coverage.overall} & \metric{erl.text.pct.h2.gedi_night.coverage.overall} \\
\hline
\end{tabular}
\end{table}
